\chapter[研究方法（データと分析設計）]{研究方法}

本章では、本研究で用いるデータ、変数定義、新規性（Novelty）指標の構築方法、ならびに提出日をイベント日とするイベントスタディと回帰分析の設計を説明する。具体的には、まず有価証券報告書から「事業等のリスク」欄を抽出し、年度間比較が可能な形に整形する。次に、リスク記述をチャンク単位に分割し、前年記述との意味的類似度に基づいて内容更新の度合い（新規性）を定量化する。最後に、提出日周辺における市場反応を、価格反応

（CAR）に加えて注目度（出来高）および不確実性指標で多面的に測定し、新規性指標との関連を検証する。

\section{データ}

\subsection{企業・期間}

本研究は、日本の上場企業を対象とし、一定期間に提出された有価証券報告書を分析対象とする。サンプルは、提出企業のうち株価データおよび統制変数が取得可能であり、かつ前年の有価証券報告書が存在して年度間比較が可能な企業・年度に限定する。

\subsection{文書データ（有価証券報告書）}

文書データとして、有価証券報告書から「事業等のリスク」欄の本文を取得する。取得は電子開示システム（EDINET）により提供される提出書類データを用い、機械的に当該セクションを抽出したうえで、分析可能なテキストへ変換する。抽出の段階では、年度間比較の妨げとなる不要な装飾情報（例：HTML 由来のタグ、記号の過剰な連続、機械抽出のノイズ）を除去し、同一企業について前年との対応づけが可能な形で整形する。

\subsection{市場データ（株価・出来高）}

市場データとして、イベントスタディに必要な日次株価・出来高データを取得する。株価データは、提出日をイベント日としてイベント窓における異常収益を算出するために用いる。また、投資家の注目度を捉えるため、出来高を市場水準で調整した異常出来高指標を構築する。さらに、企業特性による交絡を抑えるため、規模や収益性、財務レバレッジ等の統制変数を取得し、回帰分析に組み込む。

\section{変数定義}

\subsection{説明変数：リスク開示の新規性（Novelty）}

本研究の中心的な説明変数は、有価証券報告書「事業等のリスク」欄における記述の新規性（前年からの内容更新の度合い）である。新規性は、当該セクションのテキストを一定単位（チャンク）に分割したうえで、前年テキストとの意味的類似度にもとづいて定量化する（構築方法の詳細は 2.3 で述べる）。本研究では、全体の新規性を表す指標（例：change\_topk）を基本とし、必要に応じてカテゴリ別の新規性指標も用いる。

\subsection{被説明変数（アウトカム）}

提出日をイベント日とし、イベント窓における市場反応を以下の 3 指標で測定する。

\begin{itemize}

  \item 注目度（出来高）：abvol\_mkt

\end{itemize}

市場平均出来高で調整した出来高（対数変換した出来高等）をイベント窓で集計した指標であり、情報に対する投資家の注意（attention）の高まりを捉える代理変数と位置づける。

\begin{itemize}

  \item 不確実性：vol\_abs

\end{itemize}

市場モデルに基づく異常収益（AR）の絶対値をイベント窓で合計した指標（\(\sum \lvert AR \rvert\)）であり、開示を受けた短期の不確実性・解釈のばらつきの大きさを捉える代理変数と位置づける。

\begin{itemize}

  \item 価格反応（参考）：car\_mm

\end{itemize}

市場モデルに基づく累積異常収益率（CAR）であり、開示が株価水準に与える短期影響を捉える。

\subsection{統制変数}

新規性と市場反応の関係が企業特性により交絡しないよう、以下の統制変数を用いる。

\begin{itemize}

  \item 企業規模：log\_assets

  \item 事業規模：log\_revenue

  \item 収益性：roa

  \item 財務レバレッジ：lev

\end{itemize}

加えて、マクロ環境や年次ショックを吸収するための年固定効果（年 FE）、業種特性を吸収するための業種固定効果（業種 FE）を導入する。

\section{新規性（Novelty）指標の構築}

本研究では、「事業等のリスク」欄の新規性を、前年テキストとの意味的差分として捉える。具体的には、当該セクションをチャンク単位に分割し、各チャンクについて前年チャンク集合との意味的類似度を算出する。類似度は文の意味表現を反映する埋め込みベクトルにもとづき評価し、当年チャンクが前年にどの程度“対応する記述”を持つかを測る。当年の各チャンク \(c\) について、前年チャンク集合 \(\mathcal{P}\) との類似度の最大値 \(\max_{p\in\mathcal{P}}\mathrm{sim}(c,p)\) を求め、変化度を \(\mathrm{change}(c)=1-\max_{p\in\mathcal{P}}\mathrm{sim}(c,p)\) として定義する。直観的には、前年に非常に近い記述が存在するチャンクは変化度が小さく、前年に近い対応記述が見当たらないチャンクほど変化度が大きい。これを企業・年度単位で集約することで、「当年のリスク開示が前年からどれだけ更新されたか」を表す新規性指標（例：change\_topk）を構築する。

\section{実証設計：イベントスタディと回帰分析}

本研究では、市場モデルにより期待収益率を推定し、異常収益（AR）を構成する。銘柄 \(i\) の日次収益率を \(R_{i,t}\)、市場収益率を \(R_{m,t}\) とすると、市場モデルは次式で与えられる。
\[
R_{i,t}=\alpha_i+\beta_i R_{m,t}+\varepsilon_{i,t}
\]
推定された \(\hat{\alpha}_i,\hat{\beta}_i\) を用いて異常収益を
\[
AR_{i,t}=R_{i,t}-\hat{\alpha}_i-\hat{\beta}_i R_{m,t}
\]
と定義する。イベント窓 \([\tau_1,\tau_2]\) における累積異常収益率（CAR）は
\[
CAR_{i}(\tau_1,\tau_2)=\sum_{t=\tau_1}^{\tau_2} AR_{i,t}
\]
と定義する。

\subsection{イベントの定義とイベント窓}

イベント日は、有価証券報告書の提出日（公開日）とする。イベント窓は提出日前後の短期窓を用い、当該窓内で abvol\_mkt、vol\_abs、car\_mm を算出する。

\subsection{異常収益と CAR（市場モデル）}

価格反応および不確実性指標の構築にあたり、市場モデルにより期待収益を推定し、異常収益（AR）を得る。CAR はイベント窓における AR の合計として算出する。また vol\_abs はイベント窓における \(\sum \lvert AR \rvert\) として算出する。

\subsection{回帰モデル}

新規性と市場反応の関係は、企業特性を統制した回帰分析により検証する。基本形は以下のとおりである。

\begin{itemize}

  \item 被説明変数：abvol\_mkt または vol\_abs または car\_mm

  \item 説明変数：新規性指標（例：change\_topk）

  \item 統制変数：log\_assets, log\_revenue, roa, lev

  \item 固定効果：年 FE、業種 FE（必要に応じて企業 FE）

\end{itemize}

これにより、リスク開示の新規性が（i）注目度、（ii）不確実性、（iii）価格反応のいずれに強く表れるかを比較可能にする。

\section{識別上の留意点と推定上の設定}

\subsection{識別上の留意点（解釈）}

本研究は、リスク開示の新規性と提出日周辺の市場反応の関連をイベントスタディおよび回帰分析により検証する。ただし、推定される関係は必ずしも厳密な因果効果を意味しない点に留意が必要である。第一に、リスク開示の新規性は企業の実態変化（業績悪化、事業環境の急変、リスク顕在化等）に伴って高まる可能性があり、その実態変化自体が市場反応を同時に引き起こし得る。第二に、企業は市場環境や投資家の関心を踏まえて開示内容を調整する可能性があり、新規性は外生的に決まるとは限らない。したがって本研究では、企業規模・収益性・財務構造といった観測可能な企業特性、ならびに年固定効果・業種固定効果を導入し、交絡の影響を可能な限り抑えた上で、推定結果を「提出日周辺における市場反応との統計的関連」として解釈する。

\subsection{推定上の設定（標準誤差・固定効果等）}

回帰分析では、統制変数として log\_assets, log\_revenue, roa, lev を含め、年固定効果（年 FE）および業種固定効果（業種 FE）を導入する。標準誤差は、同一企業における誤差項の相関を考慮する観点から、基本的には企業レベルでのクラスタリング等、頑健な推定を用いることが望ましい。必要に応じて企業固定効果（企業 FE）を導入し、企業に固有で時間不変な要因（開示方針、事業構造、恒常的リスク水準等）を吸収した推定も行い、結果の頑健性を確認する。

\subsection{頑健性検証（事前窓・プラセボ偽イベント日）}

本研究では、推定結果が提出日固有の反応を捉えているかを確認するため、事前窓およびプラセボ（偽イベント日）を用いた頑健性検証を行う。具体的には、（i）提出日前の事前窓（例：-3 日から-1 日）において、新規性指標がアウトカムを予測しないこと（プレトレンドが弱いこと）を確認する。加えて（ii）提出日を外生的に±10〜30 日程度シフトした偽イベント日を設定し、その偽イベント日周辺のアウトカムに対して同一の回帰仕様を適用する。真の提出日にのみ反応が集中し、偽イベント日では効果が消失することは、推定結果が一般的な時系列相関や企業の恒常的特性ではなく、提出日に紐づく反応を捉えていることを支持する。

\section{分析環境・再現性}

\subsection{使用ソフトウェア・実行環境}

本研究の分析は、主として Python（Jupyter Notebook） により実施した（実行環境の基準として Python 3.10.18 を用いた）。データの保管・結合には PostgreSQL を用い、リスク開示テキストの処理、指標構築、イベントスタディ、回帰分析、図表作成までをノートブック上で一貫して実行する。テキストの意味表現（埋め

込み）には sentence-transformers を用い、埋め込みモデルとして intfloat/multilingual-e5-small を使用した（実装上は環境変数 RISK\_EMBED\_MODEL によりモデル指定を変更可能）。これらの埋め込みは BERT 系モデルに基づく文表現の枠組みの上に位置づけられ\cite{Devlin2019,ReimersGurevych2019,Wang2022E5}、本研究では日本語テキストに対して多言語埋め込みモデルを適用した。チャンク間の類似度は、埋め込みベクトルにもとづくコサイン類似度により算出する。

\subsection{コード公開と参照方法}

本研究で用いた分析コード（ノートブック、前処理・推定・図表作成コード）は GitHub リポジトリとして公開し、第三者が参照・再実行できる形で提供する。リポジトリのURLは次のとおりである。
\begin{quote}
\url{https://github.com/Junxzi/yuho-risk-novelty-analysis}
\end{quote}
実行手順、依存関係、必要な環境変数、ならびに生成物の保存先はリポジトリ内の README に記載する。なお、取得データ（有価証券報告書本文テキスト、株価・財務データ等）は利用条件に配慮し、リポジトリには同梱しない。

\subsection{設定・乱数・キャッシュ}

分析の計算効率と再現性を確保するため、中間生成物や推定用データセットをファイルとして保存し、以降の処理で読み込む（キャッシュ）運用を行う。キャッシュや推定結果、図表は主として data/processed/ および outputs/ 配下に出力される。サンプル抽出等で乱数を用いる場合は seed を固定し、ノートブック内の設定（例：SAMPLE\_MODE、SAMPLE\_SEED 等）を明示したうえで実行する。

\subsection{データ入手可能性と再現上の注意}

有価証券報告書データは公開情報に基づくが、株価・出来高や一部の財務データは取得元の利用規約・API 制約に依存する場合がある。第三者が再現する際には、同一の取得元・期間・ユニバース条件を満たすデータを準備し、README に示す手順に従って処理を行う必要がある。また、提出日（イベント日）の定義やイベント窓設定は結果に影響し得るため、本研究で用いる設定値はコード上で一元管理し、変更した場合は当該箇所を明示する。
\medskip
再現性の観点から、Python のバージョンおよび主要ライブラリのバージョンは \texttt{requirements.txt}（リポジトリ直下）に記載する。第三者が実行環境を再構築する際には、当該ファイルに従って依存関係を揃えることが望ましい。
